\begin{SCn}
\begin{small}

\scnheader{CycL}
\scnidtf{Язык описания онтологий, используемый в проекте искусственного интеллекта Cyc}
\scnidtf{CycL (Cyc Language)}
\scniselement{язык описания онтологий}
\begin{scnrelfromlist}{разработчик}
    \scnitem{Дуглас Ленат}
    \scnitem{Раманатан В. Гуха}
\end{scnrelfromlist}
\scnrelfrom{год создания}{1984}
\scnrelfrom{основная парадигма}{декларативное программирование}
\scntext{описание}{CycL --- формальный язык представления знаний, разработанный для проекта Cyc, одной из самых масштабных в мире баз знаний, целью которой является кодирование здравого смысла и сложных концепций через формализацию общих знаний. Программы на CycL состоят из утверждений о мире, включающих аксиомы, правила вывода и объекты, а выполнение осуществляется через алгоритмы автоматического доказательства теорем, учитывающие семантическую связанность и иерархию понятий.}
\scntext{описание}{CycL поддерживает формализацию отношений между объектами через аксиомы, правила и утверждения, позволяя строить глубокие логические выводы. Обеспечивает динамическое обновление баз знаний. Язык активно используется для создания экспертных систем, анализа данных и решения задач, требующих контекстного понимания.}
\begin{scnrelfromset}{основные конструкции}
    \scnfileitem{факт}
    \scnfileitem{правило}
    \scnfileitem{микротеория}
\end{scnrelfromset}
\scnrelfrom{стратегия вывода}{Логический вывод с учетом контекста}
\begin{scnrelfromlist}{поддержка}
    \scnfileitem{контекстное программирование}
    \scnfileitem{онтологическая модульность}
    \scnfileitem{встроенные правила наследования}
\end{scnrelfromlist}
\begin{scnrelfromlist}{применение}
    \scnitem{искусственный интеллект}
    \scnitem{представление знаний}
    \scnitem{реализация систем, обладающих здравым смыслом}
    \scnitem{обработка естественного языка}
    \scnitem{машинное обучение}
    \scnitem{сложный логический вывод}
\end{scnrelfromlist}
\begin{scnrelfromset}{примеры использования}
    \scnfileitem{Создание баз знаний, связывающих симптомы, заболевания и методы лечения, с возможностью логического вывода на основе здравого смысла}
    \scnfileitem{Реализация чат-ботов или виртуальных помощников, способных понимать контекст запросов и давать логически обоснованные ответы}
    \scnfileitem{Анализ тональности и выявление скрытых смыслов в документах}
    \scnfileitem{Аннотирование данных, что улучшает обучение моделей в задачах вроде распознавания образов или прогнозирования поведения пользователей}
    \scnfileitem{Создание интеллектуальных тренажёров, которые анализируют ошибки ученика и предлагают персонализированные пути обучения, опираясь на здравый смысл}
\end{scnrelfromset}
\begin{scnrelfromset}{реализации}
    \scnitem{Cyc Core}
    \scnitem{OpenCyc}
\end{scnrelfromset}
\begin{scnrelfromset}{особенности}
    \scnfileitem{строгая типизация через микротеории}
    \scnfileitem{высокая выразительность}
    \scnfileitem{интеграция с онтологиями}
    \scnfileitem{поддержка неполных знаний}
\end{scnrelfromset}
\begin{scnrelfromset}{ограничения}
    \scnfileitem{специфичность для системы Cyc}
    \scnfileitem{ограниченная гибкость вне контекста онтологий}
\end{scnrelfromset}
\begin{scnrelfromset}{библиографические источники}
    \scnfileitem{Lenat D. B., Guha R. V. Building Large Knowledge-Based Systems. Addison-Wesley, 1990.}
    \scnfileitem{Lenat, D.B. The Dimensions of Context Space / D.B. Lenat, R.V. Guha. — Stanford, CA: Stanford University, 1989. — 152 p.}
    \scnfileitem{Matuszek, C. Using Cyc to Enable Natural Language Interaction with a Commercial Game / C. Matuszek [и др.] // Proceedings of the 20th International Florida Artificial Intelligence Research Society Conference (FLAIRS-20). — Menlo Park, CA: AAAI Press, 2007. — P. 456–461.}
    \scnfileitem{CycL // Wikipedia – 2025. – Электронный ресурс. Дата обращения: 24.05.2025. https://en.wikipedia.org/wiki/CycL}
    \scnfileitem{Документация Cyc – 2025. – Электронный ресурс. Дата обращения: 25.05.2025. https://www.cyc.com}
\end{scnrelfromset}

\scnheader{LISP}
\scnidtf{Функциональный язык программирования, разработанный Джоном Маккарти}
\scnidtf{LISP (LISt Processing language)}
\scniselement{функциональный язык}
\begin{scnrelfromlist}{разработчик}
    \scnitem{Джон Маккарти}
\end{scnrelfromlist}
\scnrelfrom{год создания}{1958}
\scnrelfrom{основная парадигма}{функциональное программирование}
\scntext{описание}{LISP --- универсальный язык функционального программирования для работы с символьными данными и списковыми структурами. Язык основан на концепции S-выражений, где код и данные представлены в единой форме, что позволяет использовать программы как данные для метапрограммирования через макросы. LISP ориентирован на обработку иерархических данных, рекурсивных вычислений и динамических структур.}
\scntext{описание}{LISP оптимален для работы с высокоабстрактными структурами данных, реализации сложных алгоритмов искусственного интеллекта и метапрограммирования. Благодаря своей гибкости, основанной на унифицированном представлении кода и данных через списки, а также мощной системе макросов, он обеспечивает высокую выразительность и адаптивность.}
\scntext{назначение}{LISP обеспеченивает гибкость и выразительность при разработке программ, особенно в областях, где требуется манипуляция списками, символьными выражениями и динамическая генерация кода.}
\begin{scnrelfromset}{основные конструкции}
    \scnfileitem{S-выражение}
    \scnfileitem{функция}
    \scnfileitem{макрос}
\end{scnrelfromset}
\scnrelfrom{стратегия вывода}{Интерпретация и компиляция функциональных выражений}
\begin{scnrelfromlist}{поддержка}
    \scnfileitem{рекурсия}
    \scnfileitem{динамическое управление памятью}
    \scnfileitem{макросы для расширения языка}
\end{scnrelfromlist}
\begin{scnrelfromlist}{применение}
    \scnitem{искусственный интеллект}
    \scnitem{обработка естественного языка}
    \scnitem{математические вычисления}
    \scnitem{формальные доказательства}
    \scnitem{прототипирование алгоритмов}
\end{scnrelfromlist}
\begin{scnrelfromset}{примеры использования}
    \scnfileitem{Парсеры для разбора грамматических конструкций или системы машинного перевода}
    \scnfileitem{Реализация систем, упрощающих математические выражения или находящих аналитические решения уравнений}
    \scnfileitem{Эксперименты с новыми алгоритмами, такими как генетические программирование или нейронные сети}
    \scnfileitem{Реализация алгоритмов для задач, требующих анализа сложных структур данных}
\end{scnrelfromset}
\begin{scnrelfromset}{реализации}
    \scnitem{Allegro Common Lisp}
    \scnitem{CLISP}
    \scnitem{Clozure CL}
    \scnitem{SBCL (Steel Bank Common Lisp)}
    \scnitem{CMUCL}
    \scnitem{ECL (Embeddable Common Lisp)}
\end{scnrelfromset}
\begin{scnrelfromset}{особенности}
    \scnfileitem{минималистичный синтаксис}
    \scnfileitem{гибкость метапрограммирования}
    \scnfileitem{формализация онтологии функций}
\end{scnrelfromset}
\begin{scnrelfromset}{ограничения}
    \scnfileitem{ограничение алфавита и регистронезависимость}
    \scnfileitem{проблемы интеграции}
    \scnfileitem{специфичность применения}
\end{scnrelfromset}
\begin{scnrelfromset}{библиографические источники}
    \scnfileitem{LISP | Artificial Intelligence, Machine Learning \& Programming // Massachusetts Institute of Technology (MIT). – Электронный ресурс. Дата обращения: 25.05.2025. https://www.britannica.com/technology/LISP-computer-language}
    \scnfileitem{Лисп // Википедия – 2025. – Электронный ресурс. Дата обращения: 25.05.2025. https://ru.wikipedia.org/wiki/Лисп}
    \scnfileitem{What is the Lisp (List Processing) Programming Language? // TechTarget. – 2022. – Электронный ресурс. Дата обращения: 25.05.2025. https://www.techtarget.com/whatis/definition/LISP-list-processing}
    \scnfileitem{Введение в программирование систем искусственного интеллекта [PDF] // Джон Маккарти. – 1958.}
    \scnfileitem{Введение в язык Лисп и функциональное программирование. // Хабр. – 2023. – Электронный ресурс. Дата обращения: 25.05.2025. https://habr.com/ru/articles/764594/}
\end{scnrelfromset}

\bigskip

\scnheader{Сравнение языка описания онтологий CycL и языка функционального программирования LISP}
\begin{scneqtovector}
    \scnitem{CycL}
    \scnitem{LISP}
\end{scneqtovector}
\begin{scnrelfromset}{сходства}
    \scnfileitem{Оба применяются в задачах искусственного интеллекта и обработки знаний.}
    \scnfileitem{Поддерживают формализацию сложных структур данных.}
    \scnfileitem{Имеют открытые и коммерческие реализации, используются в научных и промышленных проектах}
\end{scnrelfromset}
\begin{scnrelfromset}{различия}
    \scnfileitem{CycL — декларативный язык для онтологий, LISP — функциональный язык общего назначения}
    \scnfileitem{CycL использует микротеории для строгой типизации, LISP — S-выражения и макросы для гибкости}
    \scnfileitem{CycL специфичен для системы Cyc, LISP универсален и применяется в широком спектре задач}
    \scnfileitem{CycL гарантирует семантическую связанность, LISP — динамическую генерацию кода}
\end{scnrelfromset}
\begin{scnrelfromset}{рекомендации по выбору}
    \scnfileitem{CycL рекомендуется для задач, требующих формализации знаний в рамках онтологий и логического вывода с учетом контекста}
    \scnfileitem{CycL подходит для проектов, связанных с обработкой естественного языка и сложными логическими выводами}
    \scnfileitem{LISP целесообразно использовать для разработки алгоритмов искусственного интеллекта, прототипирования и метапрограммирования}
    \scnfileitem{LISP оптимален для задач, где важна гибкость, динамическое управление памятью и поддержка рекурсии}
\end{scnrelfromset}

\end{small}
\end{SCn}
